\clearpage\SourceLeaf{293}

\begin{center}
\textbf{BEMERKUNGEN ZU DEN NOTIZEN ÜBER GEOMETRIA SITUS.}
\end{center}

Die unter [I] zusammengefassten Notizen stehen hinter einander in einem
Handbuche, sind jedoch augenscheinlich zu sehr verschiedenen Zeiten
niedergeschrieben worden, die Notizen [1] und [2] stammen wahrscheinlich
aus der Zeit zwischen 1823 und 1827, während die Notiz [9] vermuthlich in
die Zeit nach 1840 fällt.

In \textsc{Scherings} Nachlass haben sich folgende Erläuterungen zu den Notizen
[I] gefunden:

»Um die Charactere $\gamma$, $\gamma'$ [S.\,272] irgend eines bestimmten
Theils der Curve zu erhalten, ziehe man von beliebig weiter Ferne nach
zwei Punkten, die unmittelbar neben dem Curventheil und auf verschiedenen
Seiten des Zuges liegen, zwei gerade Linien und zähle für jeden
Durchschnitt, welchen eine der beiden Linien mit dem gegebenen Zuge
macht, $+1$, wenn die Linie von der äussern Seite zur innern übergeht,
$-1$ im entgegengesetzten Falle. Die Resultate für die beiden Linien sind
die gesuchten Charactere $\gamma$, $\gamma'$.

Nennt man $+1$ oder $-1$ den Character eines durch eine einzelne der in
sich zurückkehrenden Linien begrenzten Flächentheils, je nachdem die
Fläche auf der positiven oder negativen Seite ihrer Begrenzung liegt, so
bezeichnet $\mathrm{Ch}\,\varepsilon a$ [S.\,274] die Summe der
Charactere aller der Flächentheile, die an $a$ stossen, sowohl derjenigen,
die über $a$ hinausgehen, als auch des Flächentheils, der von
$\varepsilon a$ begrenzt wird und dessen Character $\varepsilon$ ist.

Der absolute Character $\chi$ eines Punktes [S.\,278] ist die Summe der
Charactere aller Flächentheile, in denen der Punkt liegt.

Bei der Zählung [S.\,281] hat man die Seiten des Polygons in einem und
demselben Sinne, immer in der Richtung von $(x, y)$ nach $(x', y')$, zu
durchlaufen.«

Die Notiz [II] hat \textsc{Schering}, wahrscheinlich auf Grund der Äusserungen
von \textsc{Gauss} in der Anmerkung S.\,272, mit dem Datum: 1844.\ Dec.\ 30
versehen und zu ihr folgende Bemerkungen gemacht:

»Zwei Punkte des Tractes, die zusammenfallen und dadurch einen Knoten
bilden, sind mit denselben Buchstaben \emph{a} und \emph{a}, \emph{b} und
\emph{b} u.\,s.\,w.\ bezeichnet. Bei der Darstellung eines Tractes werden
diese Punkte in der Ordnung angegeben, in welcher sie auf dem Tracte nach
einander folgen.

Durch die Figuren werden die Nummern

\begin{center}
1 \quad 54 \quad 15 \quad 3 \quad 10

15
\end{center}

dargestellt.«\hfill \textsc{Stäckel}.
